\chapter{Выбор необычной архитектуры после доказательства спектральной слепоты}

Обычные спектральные данные сохраняют сумму матриц Грама в средней вершине,
но забывают её ориентированное разложение. Следующая архитектура обязана
содержать данные, чувствительные к ориентации. Здесь сравниваются оставшиеся
классы и выбирается одна наиболее дешёвая точная проверка.

\section{Условия допустимости}

Кандидат допускается к сравнению, только если он не является обычной
функцией от \(\Spec D^2\), не совпадает с уже закрытой реализацией,
порождает структуру «входящее минус исходящее» до обращения к феноменологии,
не вводит подбираемый непрерывный коэффициент и допускает конечную проверку
на существующей поколенческой цепочке KO6.

\section{Сводная таблица направлений}

\begin{center}\small
\begin{tabular}{>{\raggedright\arraybackslash}p{3.1cm}
                  >{\raggedright\arraybackslash}p{2.2cm}
                  >{\raggedright\arraybackslash}p{2.2cm}
                  >{\raggedright\arraybackslash}p{4.6cm}}
\toprule
Класс & Состояние & Точная цель & Решение\\
\midrule
Кривизна конуса отображения & закрыт & нет & сохраняет произведение концов и теряет чистые квартичные члены\\
Вещественное вспомогательное поле & закрыт & нет & исключение поля даёт неверный знак\\
Комплексное преобразование Хаббарда--Стратоновича & мера KO6 закрыта & лишь формально & требует комплексной седловой точки\\
Скрученное спектральное исчисление & новый класс & пока нет & допустимые флуктуации расширяются, но требуемое отображение момента не выведено\\
Производная симплектическая модель и BV--BFV & новый класс & возможно & язык редукции подходит, но конечный комплекс проекта отсутствует\\
Относительный модулярный функционал & открыт слабо & нет & асимметрия есть, но локальный гессиан сам по себе не выбирает разность\\
Нелокальное граничное действие & текущая реализация закрыта & нет & требуется новое гильбертово пространство и новый след\\
Ориентированное действие высоты--Ходжа & не проверен & да & использует существующую цепочку, целую высоту и средний след\\
\bottomrule
\end{tabular}
\end{center}

Скрученные тройки действительно расширяют исчисление флуктуаций
(\path{arXiv:1411.1320}); производная пуассонова редукция и BV--BFV дают
симплектический язык редукции (\path{arXiv:2012.04451},
\path{arXiv:2106.06625}, \path{arXiv:1905.08047}); отображения момента
колчанов естественно возникают в некоммутативной симплектической геометрии
(\path{arXiv:math/0502301}). Однако ни одна из этих работ не вычисляет
конкретный конечный поколенческий блок проекта.

\section{Извлечение ориентированного дифференциала из высоты}

На трёхвершинной цепочке положим
\begin{equation}
 h=\operatorname{diag}(-I_3,0_3,I_3),\qquad D=d+d^\dagger.
\end{equation}
Две направленные стрелки повышают высоту на единицу, поэтому
\begin{equation}
 [h,d]=d,\qquad[h,d^\dagger]=-d^\dagger.
\end{equation}
и
\begin{equation}
 \boxed{d=\frac12\bigl(D+[h,D]\bigr).}
 \label{eq:v5-height-oriented-differential}
\end{equation}
Обращение \(h\mapsto-h\) меняет местами \(d\) и \(d^\dagger\), но не
изменяет квадрат их коммутатора Ходжа.

Для
\begin{equation}
 d=\begin{pmatrix}0&0&0\\X&0&0\\0&Y&0\end{pmatrix}
\end{equation}
получаем
\begin{equation}
 [d,d^\dagger]\big|_{V_G}=XX^\dagger-Y^\dagger Y.
 \label{eq:v5-height-middle-moment-map}
\end{equation}

\section{Средний проектор как многочлен от высоты}

Единственную вершину нулевой высоты выделяет проектор
\begin{equation}
 \boxed{P_G=I-h^2=\operatorname{diag}(0_3,I_3,0_3).}
 \label{eq:v5-height-middle-projector}
\end{equation}
Среди аффинных многочленов от \(h^2\) условия \(q(0)=1\), \(q(1)=0\)
однозначно дают \(q(h^2)=I-h^2\).

Определим
\begin{equation}
 S_{hH}(D,h)=\frac13\Tr_{\mathcal H}
 \left(P_G[d,d^\dagger]^2\right),
 \qquad d=\frac12(D+[h,D]).
 \label{eq:v5-height-hodge-action}
\end{equation}
Тогда точно
\begin{equation}
 \boxed{S_{hH}=\tau_3\left((XX^\dagger-Y^\dagger Y)^2\right).}
 \label{eq:v5-height-hodge-exact-target}
\end{equation}
При \(Y=\Phi I_3\) это требуемое выражение
\(\tau_3((XX^T-|\Phi|^2I_3)^2)\). Блочно-диагональные калибровочные
преобразования коммутируют с \(h\) и \(P_G\), а положительность следует из
самосопряжённости \([d,d^\dagger]\).

\section{Отличие от неудачной конструкции конуса отображения}

Предыдущая относительная конструкция проецировала квадрат
\((d+d^\dagger)^2\) и сохраняла произведение концов пути. Здесь исходной
кривизной служит уже блочно-диагональный ориентированный коммутатор
\([d,d^\dagger]\), после чего \(I-h^2\) оставляет его среднюю компоненту.
Это другой тип кривизны и другой функционал.

\section{Нерешённый вопрос происхождения}

Алгебраическая формула точна, но необходимо вывести высоту из алгебры,
градуировки и вещественной структуры KO6 с точностью до общего знака;
согласовать её с частичной и античастичной цепочками; получить из одного
нормированного следа кинетические члены и \(S_{hH}\); установить
геометрический смысл действия; повторить полный гессиан и проверку одной
моды Боголюбова--де Жена.

\begin{nogocriterion}[Скрытый выбор высоты]
Если полные условия KO6 допускают две неэквивалентные высоты или независимые
относительные веса следа, недостающий селектор лишь спрятан в \(h\), и
конструкция должна быть отвергнута.
\end{nogocriterion}

\section{Вывод}

\begin{protocol}
\begin{enumerate}
 \item обычное спектральное, конусное, вещественное вспомогательное и
 гауссово направления: \textbf{закрыты};
 \item скрученное и производное направления: \textbf{открыты, но отложены};
 \item точное действие высоты--Ходжа: \textbf{получено};
 \item единственное происхождение \(h\), \(P_G\) и следа из KO6:
 \textbf{не установлено};
 \item физическое замыкание: \textbf{не достигнуто}.
\end{enumerate}
\end{protocol}

Следующая проверка --- \path{version5_oriented_height_hodge_ko6_gate}.
Машинный реестр сохранён в
\path{s2t_v5_nonordinary_architecture_fork_gate_results.json}.